% !TEX TS-program = lualatex+se+debug
% !TEX encoding = UTF-8 Unicode

\documentclass[14pt, twoside, a4paper]{extreport}
% -----------------------------------------------------------------------------------
% ------------------------- USTAWIENIA STRONY--------------------------
% -----------------------------------------------------------------------------------
\usepackage{geometry}
\geometry{
	left=10mm,
	right=10mm,
	top=15mm,
	bottom=15mm,
	marginparwidth=0mm,
	marginparsep=0mm
	}
\newcommand{\wstawplik}[1]{\input{../include/#1}}%skraca kod
\wstawplik{headerofficium}
% elementy zmienne co do headerofficium:
\graphicspath{{../img/}} % katalog obrazów
%\svgpath{{../include/}}
\grechangestaffsize{20}%rozmiar wyniku GregorioTEX
% koniec 
\gresetlastline{ragged}% justowanie do lewej
% debugowanie:
% %\iffalse %*** %\fi %***
% \tracingall przed problematycznym kodem - bardzo dużo informacji
% lualatex+se+debug - dodatkowe opcje: -file-line-error i -halt-on-error 
% \errorcontextlines
%\def\mydebug{1}
\begin{document}
\makeatletter
\grechangedim{interwordspacetext}{\unexpanded{\the\fontdimen2\font plus \the\gluestretch\gre@space@skip@interwordspacenotes minus \the\glueshrink\gre@space@skip@interwordspacenotes}}{fixed}%
\makeatother
\pagestyle{empty}
%\pagestyle{plain}	
\large
\begin{spacing}{0.7}
\begin{center}
	\section*{Bezpieczeństwo żywnościowe Polaków}
	%\small\textit{Fontevraud Gradual\\}
	\small\textit{prof. Grażyna Cichosz (SEJM RP - 29.04.2025)}
	\vspace{4mm}
\end{center}	
\end{spacing}
\setstretch{1}% odstęp między liniami
\small
\selectlanguage{polish}
%\noindent\subsekcja{ROZDZIAŁ I.}
(0:00) Może jeszcze dodam, że mój staż z zakresu technologii to jest 52 lata, (0:05) staż z zakresu bezpieczeństwa żywności 23 lat, czyli jest trochę wiedzy. (0:12) Szanowni Państwo, jestem na finiszu z opracowaniem książki (0:15) pod tytułem Kontrowersje wokół żywności. (0:18) Więc siłą rzeczy muszę być na bieżąco z literaturą.
(0:22) Informuję wszystkich z pełną odpowiedzialnością, (0:26) że bezpieczeństwo zdrowotne żywności zapisane w prawie żywnościowym (0:30) czy w zaleceniach żywieniowych jest totalną iluzją. (0:36) Zasada przezorności, która rzekomo obowiązuje w rolnictwie, (0:41) jeżeli coś może być szkodliwe, to nie jest wdrażane. (0:45) Zasada ostrożności, która rzekomo obowiązuje w produkcji żywności, (0:52) również nie obowiązuje, to są puste zapisy, nic nieznaczące.
(1:00) Dałam tekst książki, który musi być przeredagowany, (1:03) skrócony generalnie i odnaukowiony, bo jest za trudny do zrozumienia, (1:08) do przeczytania specjalistom, no i doskonale moje intencje namierzyli. (1:14) Piszesz o tym, że to jest specjalne, celowe działanie, bo jest. (1:19) Zapewniam Państwa, jest to świadome, celowe działanie.
(1:23) To nie jest kwestia przypadku. (1:26) To nie jest kwestia przypadku. (1:29) Nie działa.
(1:34) Już jest nie włączone. (1:35) Dobra, spokojnie. (1:37) Pan tutaj będzie poprawiał, a ja kontynuuję.
(1:41) Słuchajcie Państwo, Komisja Europejska tylko trzy lata pracowała (1:46) nad sposobem znakowania żywności wartością odżywczą. (1:50) W pierwszej wersji w ogóle nie uwzględniono białka, (1:55) a przecież wartość odżywcza to białko. (1:57) To aminokwasy egzogenne, ich ilość, proporcje między nimi.
(2:01) Po interwencji naszego Ministerstwa Zdrowia wprowadzono zapis o tym, (2:09) że na etykiecie ma być podana zawartość białka. (2:12) Niewiele z tego dla nas wynika, albowiem nie dowiemy się z tej informacji, (2:16) czy jest to pełnowartościowe białko zwierzęce, (2:21) czy niepełnowartościowe, ciężkostrawne w dodatku białko roślinne, (2:26) czy produkt zawiera utlenione, glikowane, nitrozowane obecne w żywności (2:31) przetworzonej aminokwasy, czy zawiera toksyczne, uwaga, wręcz niestrawne, (2:38) wręcz toksyczne białka sojowe, które są wszechobecne w naszej żywności. (2:43) Zapis o obecności owadów i robaków, z tego co mi wiadomo, (2:47) w wersji węgierskojęzycznej zastrzegli sobie Węgrzy.
(2:52) Czy nasze agendy odpowiedzialne za bezpieczeństwo to zastrzegły? (2:58) Nie wiem, ale mogę państwu powiedzieć jedną informację. (3:02) Najwięcej robaczy... (3:02) Przepraszam, pani profesor, jeszcze raz, niech pani to powtórzy, (3:05) bo ludzie niektórzy po prostu wyłączają czujność. (3:08) O tych Węgrach jeszcze raz, jakby pani mogła powtórzyć.
(3:11) Węgrzy zastrzegli, że na etykietach produktów spożywczych, (3:17) w których są te wynalazki, ma być podana informacja (3:23) z węgiersko brzmiącą nazwą owada czy robala. (3:28) A w Polsce tego nie ma. (3:29) Nie wiem, nie wiem, po prostu nie znalazłam informacji.
(3:32) Obawiam się, że nie. (3:35) Na etykiecie mamy podane ilości kwasów tłuszczowych (3:40) nasyconych i jednowielonienasyconych. (3:42) Kto wie, jakie są proporcje w olejach czy w tłuszczach zwierzęcych? (3:46) Czyli dezinformacja totalna.
(3:49) Również na etykiecie nie znajdziemy informacji, (3:52) czy stosowane są syntetyczne rakotwórcze antyoksydanty, (3:58) które zapobiegają utlenianiu tłuszczów roślinnych. (4:01) Dlaczego nie znajdziemy informacji? (4:03) Bo w momencie, kiedy pani redaktor Katarzyna Bosacka (4:07) z dr Zdrówko Świętej Pamięci prof. Małgorzatą Kozłowską-Wojciechowską (4:12) rozkręciły akcję Skracamy etykiety, (4:15) wówczas lobbyści zmienili przepisy prawa. (4:18) I jeżeli coś jest dodawane w ilości poniżej 2\%, (4:22) to tego nie ma.
(4:23) Krótko. (4:24) Dużą część dodatków przestaną znakować literką E, (4:30) rzekomo pochodzą z naturalnych źródeł, (4:33) ale jak są izolowane, ekstrachowane, to już jest temat tabu. (4:37) I dzięki temu ilość dodatków zmniejszyła się do 10 tysięcy zaledwie.
(4:43) Dzięki tym machinacjom. (4:45) Proszę następnego. (4:48) Jakość żywności zależy od jej stabilności oksydacyjnej, (4:52) czyli od zawartości białka, które działa antyoksydacyjnie.
(4:55) I tu najcenniejsze są aminokłasy siarkowe. (4:58) Od zawartości witamin, zwłaszcza rozpuszczalnych w tłuszczach. (5:01) Tymczasem w naszej żywności, dzięki wdrożeniu wysokowydajnych technologii (5:05) w produkcji zwierzęcej, gdzie żywienie tradycyjne (5:09) zielonką paswiskową, sianem, sianokiszonkami (5:11) zastąpiono paszą przemysłową na bazie kukurydzy, soli i pszenicy.
(5:18) Dzięki tym technologiom zawartość naturalnych antyoksydantów (5:22) w mięsie i w mleku zmalała 3-4 krotnie. (5:25) Natomiast zawartość tych najcenniejszych, działających prozdrowotnie, (5:31) m.in. antynowotworowo, składników takich jak skoniugowany kwas rinolowy CLA (5:36) zmalała kilkunastokrotnie. (5:39) Na superzbilansowanej diecie monitorowanej komputerowo (5:44) krowy żyją 4 lata, a na trawce dożywają 20.
(5:50) No niestety dają mniej mleka. (5:52) I bez problemu wydają potomstwo. (5:56) Więc dzięki wysokowydajnym technologiom w produkcji zwierzęcej (6:01) potencjał prozdrowotny mięsa i mleka został całkowicie zaprzepaszczony.
(6:09) Jeżeli zawartość bioaktywnych składników jest kilkunastokrotnie mniejsza, (6:13) no to i efekty prozdrowotnego działania, jakie są z tego powodu w Szwajcarii, (6:19) po zaledwie 6 latach stosowania przemysłowej hodowli bydła mlecznego (6:25) zrezygnowano z tego eksperymentu. (6:28) Szwajcarzy w referendum podjęli decyzję, że będą płacić więcej, o 20-25\%, (6:33) i chcą mieć oryginalne, zdrowe, mleczne produkty, a nie coś mlekopodobnego. (6:39) No u nas nie wiem, czy referendum by przeszło.
(6:42) Szkoda komentować. (6:44) Więc jak wspomniałam, stabilność oksydacyjna, czyli podatność żywności (6:48) na utrenianie zależy od tych składników, które z paszy przechodzą do mleka, (6:53) przechodzą do mięsa. (6:55) Witaminy C, witamina E, witamina A, beta-karoten.
(7:00) Enzymy antyoksydacyjne z grupy oksydoreduktaz, nie będę ich wymieniać, (7:04) skoniugowany kwas linolowy, fosfolipidy unieczyniają tlen singletowy. (7:09) To jest cząsteczka tlenu, która na ostatniej orbicie nie ma jednego elektronu (7:13) i jest przez to 1560 razy bardziej aktywna oksydacyjnie od cząsteczki tlenu, (7:21) który ma komplet elektronów. (7:23) Składnikiem, który w pierwszej kolejności ulega oksydacji, (7:28) są nienasycone kwasy tłuszczowe, bo taka jest ich budowa.
(7:33) I te nienasycone kwasy tłuszczowe utleniają po kolei wszystkie składniki, (7:38) które są obecne w żywności, przede wszystkim aminokwasy siarkowe. (7:42) Aminokwasy siarkowe decydują o homeostazie, o równowadze proiantyoksydacyjnej w naszym organizmie, (7:50) ale muszą być niezdenaturowane, nieutlenione. (7:53) Podobnie działa elchistydyna, zawierający elchistydynę dipeptyd karnozyna, cynkselen.
(8:00) Te składniki są obecne w zasadzie tylko w żywności pochodzenia zwierzęcego. (8:04) W zielskach niektórych są obecne, ale w formie niebiodostępnej, (8:09) czyli są, a jakoby ich nie było. (8:11) W żywności, która nie zawiera odpowiedniej ilości antyoksydantów dochodzi do utleniania albo glikacji, (8:20) czyli łączenia białek z cukrami, przez co wartość odżywcza spada drastycznie.
(8:26) I teraz podam Państwu przykłady, dla których lekarze wierzą święcie. (8:33) Po konferencji w Katowicach 8 marca dostałam kilkanaście maili. (8:37) Czy czerwone mięso jest trucizną? (8:39) Może być trucizną i oto dowód.
(8:43) W świeżym mięsie wołowym z hodowli przemysłowej (8:46) zawartość niestrawnych, niebiodostępnych, skarbonylowanych pochodnych (8:53) jest 24 razy większa niż w mięsie owczym pozyskanym od zwierząt żywionych tradycyjnie. (9:02) Aktualnie zaledwie 9\% światowej produkcji wołowiny (9:05) i około 30\% światowej produkcji mięsa, owiec i kus (9:09) zapewniają systemy wypasu głównie na obszarach górskich. (9:13) Skutkiem przemysłowej hodowli zwierząt (9:15) jest zwiększona podatność mięsa przetworów mięsnych, (9:19) a także mleka na utlenianie.
(9:24) Żywność jest stabilna oksydacyjnie, o czym wszyscy producenci tejże wiedzą kilka dni. (9:31) Wołowina wytrzymuje pięć dni. (9:34) W zasadzie już czwartego dnia zaczyna się przebarwiać, traci kolor (9:37) i jest sygnał, że coś się dzieje.
(9:40) Jeżeli poleży dzień dłużej, to aromatem daje znać, że się zepsuła. (9:45) No więc przemysł spożywczy robi wszystko, żeby temu zapobiec. (9:49) Ale temu się nie da zapobiec.
(9:51) Można co najwyżej zamaskować, zatuszować, oszukać konsumenta. (9:57) Na slajdzie są podane trzy badania (10:00) trzech różnych zespołów naukowych z trzech różnych krajów. (10:04) Podczas dziesięciu dni dojrzewania postmortem (10:07) utlenianiu ulega do 35\% białek miofibrylarnych.
(10:12) W pakowanej próżniowo wołowinie po 16 dniach przechowywania (10:17) stwierdzono dwukrotny wzrost stężenia pochodnych karbonylowych, czyli niebiodostępnych. (10:24) Również modyfikowana atmosfera, niezależnie czy tlen mieszany z azotem (10:29) czy z utlenkiem węgla powoduje wzrost stężenia pochodnych karbonylowych (10:34) do 30\% po 5-10 dniach, czyli ani pakowanie próżniowe, (10:42) ani modyfikowana atmosfera nie zabezpieczają wołowiny (10:46) przed utlenianiem i karbonylacją białek, czyli przed utratą wartości odżywczej. (10:54) Przechowywanie mięsa to drastyczny spadek wartości odżywczej (10:59) już po 5 dniach i teraz proszę nie pospadać z krzeseł, (11:04) jak to się ma do przepisów tak zwanego prawa żywnościowego.
(11:10) Więc mimo to, zgodnie z wytycznymi Europejskiego Urzędu do Spraw Bezpieczeństwa Żywności (11:15) w 2017 roku okres przechowywania w temperaturze 3 do 8 stopni (11:20) świeżego mięsa pakowanego próżniowo, bądź w zmodyfikowanej atmosferze, (11:25) co zresztą nic nie daje, o czym już Państwo wiecie, określono na 10 dni. (11:31) Po 5 dniach już mamy 30\% straty wartości odżywczej. (11:35) W 2020 roku ten maksymalny okres przechowywania zwiększono do 13 dni, (11:42) a w 2023 pod presją sieci handlowych ustalono maksymalny okres trwałości (11:49) dla wołowiny 23 dni, 3 tygodnie i trochę, dla jagnięciny 27, dla mięsa wieprzowego 18 dni.
(11:57) Nie znalazłam żadnych badań, jaki jest zakres utleniania oksydacji białek, (12:03) czy to jest 70\%, czy 80\%, czy 90\%. (12:08) Pretekstem do wydłużenia tego okresu przechowywania był fakt, (12:13) że właśnie w tym okresie bakteria wytwarzająca jadkie łbasiane, (12:18) czyli Crostidium botulinum jeszcze się nie rozwinęła. (12:21) Oczywiście tak zwani eksperci unijni wiedzą, że to jest przekręt, (12:28) dlatego że, uwaga, dopuścili te terminy dwudziestoparodniowe pod warunkiem, (12:37) że odpowiedzialność za bezpieczeństwo produktu ponosi producent.
(12:42) Czyli wiedzą, że jest przekręt. (12:45) Po kilku dniach przechowywania mięso nie zawiera ani enzymów o działaniu antyoksydacyjnym, (12:52) tych naturalnych, które są w mięsie, ani naturalnych przeciwutleniaczy, (12:56) które z paszy przechodzą do mleka czy do mięsa. (13:00) W trakcie przechowywania, w trakcie obróbki termicznej możliwy jest rozpad cząsteczki chemu.
(13:06) I uwolnione z cząsteczki chemu żelazo intensyfikuje te procesy, bo działa prooksydacyjnie. (13:15) I teraz moje wątpliwości technologa żywności. (13:19) Czy nastrzykiwana wołowina, bo ona jest nastrzykiwana, żeby nie straciła koloru.
(13:26) Nastrzykiwana jest specjalną solanką. (13:29) W skład tej solanki wchodzi m.in. genetycznie modyfikowana soja w ilości do 4\%, (13:35) bo przy większej ilości jest posmak grochóweczki. (13:40) Wchodzą fosforany, cukier, różne związki, które zmniejszają tzw.
aktywność wody (13:46) i zapobiegają rozwojowi mikroflory w mięsie i wydłużają trwałość. (13:52) Wszystkie te związki działają szkodliwie na nasz mikrobiom, (13:55) bo jeżeli mikroflora w żywności jest ukatrupiona, to czemu ma nie być ukatrupiona w naszym przewodzie pokarmowym? (14:02) Czy tak nastrzykiwana wołowina po 23 dniach przechowywania zawiera biodostępne, (14:09) nieutlenione, nie glikowane aminokwasy? (14:13) Wiem z całą pewnością, że nie zawiera nieutlenionych aminokwasów siarkowych, (14:18) bo taką publikację Hiszpanie wyprodukowali. (14:22) Aminokwasy siarkowe są najważniejsze dla homeostazy proantyoksydacyjnej, (14:27) głównie dla mózgu, więc ta głupawka, którą ja obserwuję na drogach, (14:32) to jest efekt żywności.
(14:34) To nie są jakieś charakteropatie, to nie są choroby natury psychicznej, (14:40) tylko to jest zapruty, niedotleniony mózg. (14:44) Zresztą jeszcze kilka innych przykładów na tę tezę państwu przedstawię. (14:48) Oczywiście wołowina przechowywana 23 dni nie zawiera żadnych witamin, (14:54) a bez witamin z grupy B również nasze komórki nie dają rady.
(14:58) Nie ma szans. (14:59) Proszę następny. (15:02) Kolejny przekręt.
(15:05) Ponieważ białka są najdroższym składnikiem żywności i co gorsza sprzyjają zdrowiu, (15:11) więc są powszechnie zastępowane hydrokoloidami. (15:15) To są tak zwane wysokowydajne technologie w przetwórstwie żywności. (15:19) Hydrokoloidy to są najczęściej wielocukry roślinne albo białka, na przykład sojowe, (15:25) które wiążą wodę.
(15:26) 1 gram hydrokoloidu, 10 gram wody. (15:30) Formalnie rzecz biorąc, większość hydrokoloidów jest dla zdrowia obojętna. (15:35) Tu wyjątkiem jest karagen, który po hydrolizie w przewodzie pokarmowym działa prozapalnie (15:42) i może generować nowotwory.
(15:46) I wyjątkiem jest genetycznie modyfikowana soja. (15:48) Skutkiem stosowania hydrokoloidów, które jak już wspomniałam same, nie są specjalnie szkodliwe. (15:54) Jest znaczny spadek zawartości białka w produktach, które tradycyjnie uznawane są za wysokobiałkowe.
(16:01) W mięsie, przetworach mięsnych, w drobiu i rybach zawartość białka może być nawet o połowę mniejsza. (16:11) Nawet o połowę. (16:13) W diecie Polaków w 2018 roku z białka pochodziło tylko 13\% energii.
(16:20) To jest zaledwie połowa tego, co powinniśmy zjadać, żeby nasz mózg i cały organizm funkcjonowały prawidłowo. (16:29) Mózg i układ odpornościowy w pierwszej kolejności są narażone na szwami przy niedoborach białka. (16:36) Jednocześnie produkty wysokobiałkowe są coraz trudniej dostępne.
(16:39) Póki co to są jaja, sery dojrzewające, bo ich się nie da zafałszować, a były takowe granty proponowane. (16:48) Wszystkie padły, zresztą mam w tym duży udział. (16:52) No i twarogi, ale uwaga, nie wszystkie twarogi, tylko te tradycyjne, które kupujemy na wagę, tak zwana krajanka.
(17:00) To są jedyne wysokobiałkowe produkty. (17:03) A kto o tym wie? (17:04) Mam nadzieję, że będziecie Państwo tą wiedzę przekazywać i upoważniam Pana do rozesłania wszystkim tego referatu. (17:12) Wszystkim bez wyjątku.
(17:13) Bardzo dziękuję, to jest bardzo ważna wiadomość. (17:16) Państwo to dostaniecie. (17:17) Tak, do popularyzacji, bo po prostu ludzie nie są świadomi.
(17:24) Zawartość białka jest dwukrotnie mniejsza, stąd i taka rozcieńczona żywność nie ma smaku, nie ma aromatu. (17:33) No więc ładuje się w nią kilogramami barwniki, aromaty, wzmacniacze smaku, konserwanty. (17:40) Około 20\% białka, jest to białko niepełnowartościowe, albo genetycznie modyfikowana soja, (17:48) albo tak zwany MOM, czyli mięso oddzielone mechanicznie, albo jedno z drugim wymieszane.
(17:53) Co się nawinie. (17:55) Dostępne w ofercie handlowej mięso kulinarne, to które kupujemy, nie wiem, na kotlety, czy na wagę. (18:02) Również zawiera hydrokoloidy, a przecież nikt przy zdrowych zmysłach nie będzie tego traktował jako żywność przetworzoną.
(18:10) Bo to jest żywność naturalna, wygląda na naturalną, ale nie jest naturalna. (18:14) Więc hydrokoloidy, barwniki, fosforany, błonnik. (18:18) De facto jest to dwu wodzian mięsa, dokładnie na kilogram mięsa serwuje się około 200 gram tych wynalazków i 800 gram wody.
(18:32) Oczywiście te receptury są utajnione, ale ja przez 10 lat miałam zajęcia na niestacjonarnych studiach w Warszawie (18:39) i kurs magisterski, tudzież doktorancki w Olsztynie, studenci sobie sami wybierali tematy seminariów. (18:46) I naściągali mi tych receptur całe multum w dwóch laboratoriach, w dwóch największych koncernach spożywczych przetwarzających mięso. (19:00) Szefowe laboratorium potwierdziły to.
Te dane są niedostępne, po prostu nieosiągalne. (19:07) Współnaukowców z Wielkiej Brytanii stwierdza, że jodan jest szkodliwy, a kilka lat później się wprowadza sól z jodanem (19:14) i jeszcze się żelazocianek dodaje. To jest świadome i celowe działanie, ale przebić szklany sufit ze świadomością, (19:24) a raczej brakiem świadomości młodych jest praktycznie niewykonalne, bo jeżeli ma certyfikat, to jest bezpieczne.
(19:33) Wrócimy do tego wątku. (19:34) To jest inteligentnie napisane E536, to jest żelazocianek. (19:41) Więc tak, dodatek soi drastycznie obniża wartość odżywczą z dwóch powodów.
(19:50) Pierwszy powód to są niepełnowartościowe białka, a drugi powód to jest obecność czynników antyżywieniowych, (19:58) które hamują wydzielanie enzymów trawiennych, przede wszystkim w trzustce. (20:05) I proszę zapamiętać, przyczyną raka trzustki wcale nie jest nadmiar słodyczy i alkoholu, (20:14) tylko dityrozyna, która występuje w żywności przetworzonej. (20:18) U braci rodzeń o tym mówiłam półtora roku temu.
(20:22) Białka sojowe nie są strawne i po to są stosowane. (20:27) Przecież trzydzieści lat temu choroby leśnioskiego Crohn'a, czy wrzodziejącego zapalenia jelit się praktycznie nie diagnozowało, (20:34) a teraz diagnozuje się u przedszkolaków. (20:39) Immunoreaktywne globuliny beta-conglicyna i glycynina nie są trawione i nie są wchłaniane w jelici cienkim.
(20:51) Białka te uszkadzają struktury przewodu pokarmowego, kosmiki jelitowe konkretnie, (20:57) a oprócz tych białek mamy w genetycznie zmodyfikowanej soli toksynę Bt i pozostałości randapu. (21:05) W konsekwencji dochodzi do przewlekłych stanów zapalnych jelita, do jego rozszczelnienia, do zaburzeń mikrobiomu, (21:14) a rozszczelnione jelito powoduje, że toksyny, antygeny białkowe wnikają do całego organizmu. (21:23) To, że mamy ponad 410 tysięcy dzieci z problemami psychicznymi, to nie jest przypadek.
(21:33) Mało tego, w soli genetycznie modyfikowanej zawartość fitoestrogenów jest pięciosiedmiokrotnie większa niż w soli tradycyjnej. (21:43) Fitoestrogeny sojowe blokują receptory estrogenowe, zaburzają procesy neuroendokrymne w życiu płodowym, (21:54) czyli już w ciąży jest człowiek narażony, przed urodzeniem człowiek jest na nie narażony, (21:59) zaburzają rozwój narządów płciowych i zachowania seksualne. (22:05) Stanowią ewidentne zagrożenie dla rozrodczości ludzi i zwierząt.
(22:09) Ja mam obietnicę Pana Bogdana, że drukujący się ostatni numer optymalny, w którym jest praca mojego autorstwa, (22:20) żywność masowego rażenia dotrze do Państwa. (22:23) I też prośba, skanujemy, wysyłamy. (22:27) Zapewniam Państwa, a współpracuję z biochemiczką medyczną z Warszawskiego Uniwersytetu Medycznego, (22:33) że gdyby genetycznie modyfikowaną soję i sól z dodatkiem jodanu i żelazocianku wyoutować z naszej diety, (22:44) to w ciągu dwóch, trzech lat zachorowalność na nowotwory tarczycy, piersi, przewodu pokarmowego (22:50) zmaleje co najmniej dwu, jeśli nie trzykrotnie.
(22:55) Co najmniej dwu, jeśli nie trzykrotnie. (22:57) Więc jak otrzymacie Państwo ten numer, skanujemy, rozsyłamy. (23:00) Proszę o powtórzenie tej wskazówki.
(23:05) Żywność masowego rażenia. (23:07) Mogę dodać krótki komentarz. (23:11) Pod rozdział w książce Rolnictwo ekologiczne, dokładnie pod tym samym tytułem, (23:17) książka Rolnictwo ekologiczne, pięć tysięcy tomów, pięć tysięcy egzemplarzy, piętnaście tomów, (23:24) wydana przez Uniwersytet Wiejsko-Miejski w Olsztynie, (23:28) została zarekwirowana przez prokuraturę w 2016 roku.
(23:34) Trzykrotnie byłam na przesłuchaniach, bo pod pretekstem, że prowadzący biuro, (23:42) grant unijny, nie zawarli na czas umowy o sporządzenie dzieła, (23:49) a standardem jest, że się umowę i rachunek wystawia od razu, (23:55) pięćdziesiąt albo siedemdziesiąt razy miałam taką sytuację, (23:58) więc pod tym pretekstem władze uczelni kazały prokuraturze zaaresztować książkę. (24:08) Rektor łgał jak znud, że musiał oddać pieniądze z grantu unijnego, (24:13) a oddał pieniądze, które przyszły, bo nikt nie dostał honorarium. (24:17) I nie, to jest najgorsze.
(24:19) Najgorsze jest to, szanowni państwo, że ze 140 autorów, (24:23) koło 40 osób to byli Europejczycy, ze wszystkich krajów europejskich, (24:28) i zresztą z Kanady też były dwie osoby, a 96 osób to były Polacy. (24:33) Z tych 96 osób, 12 osób było drhab, więc mogły się o swoje awanse obawiać. (24:41) Reszta miała tytuły, kropki i nic im nie groziło.
(24:46) Zgadnijcie, ja siedziałam tydzień czasu z mężem, z sąsiadką prawniczką. (24:53) Rektor, ówczesny minister rolnictwa, pan Gowin, (24:57) wstrzymywał rektorowi zatwierdzenie nominacji, bo wiedział, że coś się dzieje. (25:04) Więc straciłam tydzień na napisanie trzech petycji.
(25:07) Do ministra edukacji, do ministra zdrowia i do ministra rolnictwa. (25:13) Byłam bliska zawału, cztery osoby zdecydowały się to podpisać. (25:19) Ja, profesor Wiąckowski, główni winowajcy, bo my o GMO pisaliśmy, (25:25) kierownik Świętej Pamięci profesor Wiąckowski, (25:28) kierownik projektu, który przypłacił to nowotworem nerki.
(25:31) Na szczęście przeszczep się udał, facet 20 lat ode mnie młodszy funkcjonuje. (25:36) I recenzentka, która akurat w tym czasie przeszła na emeryturę. (25:41) 4 z 96.
(25:45) I to jest potwierdzenie, o czym pan poseł mówi. (25:49) Głupawka totalna, nie tylko na drogach, ale w środowiskach naukowych chyba jeszcze gorsza. (25:56) Mój mąż stwierdzi, że poziom kretynizmu jest proporcjonalny do wykształcenia.
(26:00) Szanowni państwo, więc ja od doktorantów się dowiedziałam po dwóch latach, co było grane. (26:05) Otóż pięciu panów profesorów, rektor oczywiście, (26:10) dwóch profesorów z uczelni, dwóch z Polskiej Akademii, (26:13) zaliczyło dwumiesięczny wujarz od Ameryki Południowej po Kanadę z rodzinką. (26:22) A wcześniej zauważyłam, że dwaj panowie z tego grona jeżdżą takimi samymi Mercedesami.
(26:30) Trochę dziwne. (26:32) Ale mój mąż mówi, znajdź pretekst i jedź do budynku Polskiej Akademii Nauk (26:37) i zobacz, jakie tam stoją. (26:39) Dwa kolejne.
(26:42) To była cena za przehandlowanie naszego zdrowia i naszych rolników. (26:50) I naszych rolników. (26:52) Znaczy ja miałam dwie takie propozycje.
(26:56) Zaczyna się od pół miliona plus fura za trzysta, czterysta tysięcy. (27:01) Bo tak to działa. (27:04) Koncerny na to stać.
(27:07) Po prostu. (27:07) Słuchajcie państwo, genetycznie zmodyfikowana soja stosowana jest w naszej żywności od około dwudziestu lat. (27:16) Wszędzie, gdzie to możliwe.
(27:17) Usiłowałam dla studentów w dziewiętnastym roku przed COVID-em zrobić wykaz produktów spożywczych. (27:23) Zrezygnowałam. (27:24) Trzysta sześćdziesiąt ich było.
(27:26) Teraz może być więcej. (27:28) Najnowszy wynalazek branży spożywczej, tym razem japoński. (27:32) Enzym transglutaminaza.
(27:35) Powoduje, że białka łączą się między sobą. (27:39) I dzięki temu tak zwanemu sieciowaniu wiązań kowalencyjnych te białka wiążą między i wodą. (27:49) Jeżeli w roztworze wodnym z dodatkiem transglutaminazy namoczymy tuszkę ryby, czy tuszkę kurczaka, czy pierś z indyka.
(27:59) To masa produktu wzrasta o dwadzieścia procent. (28:02) Jego trwałość oczywiście maleje. (28:06) Wartość odżywcza też maleje.
(28:08) Ale kto by się tym przejmował? (28:10) I proszę zapamiętać to, co jest wyboldowane. (28:13) Transglutaminaza jest stosowana do wykorzystywania produktów ubocznych przemysłu spożywczego. (28:21) Z tych produktów ubocznych produkuje się całkiem smaczny, atrakcyjny, sensorycznie.
(28:26) Wędlinki drobiowe, paluszki rybne, nuggetsy, wyroby garmażeryjne. (28:31) Nie byłoby problemu, bo produkty odpadowe przemysłu mięsnego są jadalne. (28:39) Problemem jest to, że wiązania, które powstają pod wpływem transglutaminazy nie są trawione w przewodzie pokarmowym, (28:47) bo nasz przewód pokarmowy się jeszcze nie nauczył ich rozpoznawać.
(28:50) Czyli w przewodzie pokarmowym znajduje się białko, które jest oporne na trawienie, (28:55) nierozpoznawalne przez układ odpornościowy. (28:57) Jeżeli to białko trafia do jelita grubego, to mikroflora w jelicie grubym robi rewolucję. (29:03) No bo ma z czego tę rewolucję robić.
(29:06) Więc wysokowydajne technologie w hodowli zwierzęcej, w przetwórstwie mięsa (29:12) to nic innego jak maksymalizacja zysków kosztem wartości odżywczej produktów półbiedy, (29:20) ale kosztem naszego zdrowia. (29:22) I tu się zaczyna problem. (29:25) Poproszę następny.
(29:29) Nasze zdrowie narażone jest na utlenione tłuszcze roślinne. (29:37) Sądzę, że państwo wiecie doskonale, że smażenie na olejach to jest prawie, że samobójstwo. (29:42) Oleje są tak zbudowane, że się nie mogą nie utleniać, a utleniają się w postępie geometrycznym.
(29:48) Kwas linolowy omega-6 utlenia się 10 razy szybciej niż kwas oleinowy z jednym wiązaniem nienasyconym. (29:55) Ten, który jest w oleju rzepakowym i w oliwie z oliwek. \\
\% do poniższych czasów dodać 29:55 \%\\
(0:06) A linolenowy omega-3 z trzema wiązaniami utlenia się 100 razy szybciej, momentalnie. (0:13) Na etapie tłuczenia się utleniają.
(0:17) Utlenione kwasy tłuszczowe działają jak wolne rodniki. (0:20) Są bardzo reaktywne i powodują utlenianie białek, zwłaszcza aminokwasów siarkowych, (0:27) tych najcenniejszych dla naszych szarych komórek. (0:30) Nadtlenki lipidowe rozkładają witaminy i te, które są rozpuszczalne w tłuszczu i witaminy z grupy B. (0:37) W konsekwencji naturalne witaminy obecne w żywności o działaniu przeciwutleniającym nie mogą być aktywne ani w produktach spożywczych, ani w naszym organizmie.
(0:48) Frytki, czipsy smażone na olejach to prosta droga do zaburzeń psychicznych i nowotworów. (0:56) Zapamiętajcie Państwo, zawierają akroleinę i akrylamid, które blokują wytwarzanie energii w mitochondriach. (1:07) A zaburzenia w funkcjonowaniu mitochondriów to jest pierwszy krok i wspólny mianownik wszystkich chorób nowotworowych i wszystkich chorób związanych z neurodegeneracją.
(1:22) To jeszcze nie wszystko. Największą fikcją i hipokryzją jest dopuszczenie dostosowania dodatków funkcjonalnych do żywności. (1:33) Aktualnie jest ich dopuszczonych około 10 tysięcy i tak formalnie rzecz biorąc każdy z nich jest przebadany.
(1:43) Każdy bez wyjątku jest przebadany. (1:45) Na zwierzętach, to są szczury też z tej samej rasy, na których się bada leki, szczury dostają coraz to większą dawkę jakiegoś tam specyfiku do momentu, kiedy połowa z nich ducha wyzionie. (1:58) I to jest wskaźnik dawka letalna dla 50\% LD50.
To LD50 dzieli się przez 100 i odnosi do wagi ciała. (2:12) Mamy inną wagę niż zwierzęta doświadczalne. (2:15) I przyjmuje się, nie sprawdzając niczego poza tym, że stosowanie tej bezpiecznej dawki czyli LD50 dzielone przez 100 w przeliczeniu na masę ciała jest bezpieczne nawet wtedy, kiedy byśmy to zjadali dzień w dzień do końca życia.
(2:34) A to nie jest prawda. (2:35) Tematem tabu pozostaje kumulacja dodatków z różnych produktów spożywczych. (2:42) Tematem tabu są interakcje między dodatkami a lekami.
(2:48) Nie uwzględnia się w ogóle biodostępności związków mineralnych, pierwiastków śladowych czy witamin. (2:53) Wszystkie napoje alkoholowe są traktowane siarczynami i witaminy z grupy B są momentalnie zablokowane. (3:01) Piwo jest doskonałym źródłem witamin z grupy B, ale one nie są wchłaniane w obecności siarczynów, więc piwo trzeba robić sobie w domu i wino również.
(3:10) W przystosowaniu i zatwierdzaniu dodatków w ogóle nie uwzględnia się wartości odżywczej, biologicznej produktu. (3:18) Co więcej, tematem tabu pozostaje wpływ dodatków na zdrowie dzieci, młodzieży, na zdrowie diabetyków, seniorów, alergików, (3:28) na zdrowie chorób z nieswoistymi stanami zapalnymi jelita, na przykład z chorobą leśnioskiego Crohna, (3:35) czy też w przypadku rekonwalescentów, na chemio, radioterapii, a nawet tylko antybiotykoterapii. (3:43) Większość dodatków jest szkodliwa dla dzieci, dla osób przyjmujących leki i szczególne zagrożenie stanowią te dodatki, (3:51) które przenikają barierę krew-mózg.
Są to aspartam, glutaminian i syntetyczne barwniki spożywcze. (4:01) Aspartam został wycofany w 2015 roku we Francji, dlatego że udowodniono jego rakotwórcze właściwości. (4:12) Na Węgrzech wykonano badania, w których wykazano, że aspartam trzykrotnie zwiększa ryzyko chłoniaka i białaczki.
(4:20) Te badania zostały zrównane z ziemią na tak zwanych konferencjach naukowych i zostały w całości powtórzone we Włoszech. (4:31) Uzyskano dokładnie takie same wyniki. U nas niestety aspartam znajduje się w sześciu tysiącach różnych produktów.
(4:40) Nie znajdziecie państwo w aptece leku dla dziecka, syropu, tabletek od bólu gardła czy innych bez aspartamu. (4:49) Pretekstem była cukrzyca. Żeby nie dawać cukru, no bo diabetyków przybywa, serwuje się aspartam.
(4:58) Przerobiłam to na własnej skórze, w siedmiu aptekach nie znalazłam leków bez aspartamu, więc moje wnuki musiały się zadowolić syropem z cebuli. (5:09) Aspartam wykazuje działanie neurotoksyczne. Powoduje zmiany ultrastrukturalne w nerwach, powoduje uszkodzenia osłonki mielinowej.
(5:19) To zwiększa ryzyko neuropatii, a to się przekłada na zwiększone ryzyko chorób nerek i wątroby, czyli osłabiony detoks organizmu. (5:27) Uwalniane z aspartamu metabolity, wszystkie bez wyjątku, są tak samo toksyczne jak on i tu szczególną neurotoksyczność wykazuje metanol. (5:37) Z jednego litra napoju sugar free, bo jeżeli jest sugar free, to nie zawiera cukru i tam jest aspartam.
(5:45) Producent nie musi podawać, że jest aspartam, bo takie jest prawo żywnościowe. (5:50) Więc z litra takiego napoju uwalnia się 60 mg metanolu. (5:57) Zdrowa dorosła osoba bez uszczerbku na zdrowiu jest w stanie zmetabolizować 7 mg metanolu.
(6:06) A ile dziecko? Nie wiem. (6:10) I ta głupawka wszechobecna, która dotyczy sądzę, że 50\% społeczeństwa albo więcej, to są przede wszystkim te dodatki, które przełamują barierę krew-mózg. (6:22) Aspartam to jest jeden z nich.
(6:24) Kolejny to jest glutaminian z sodu. (6:27) Tylko 1,5 mln ton glutaminianu trafia rocznie do żywności, oczywiście w skali globalnej. (6:34) Glutaminian stosowany jest wszędzie, gdzie to możliwe.
(6:38) Ponieważ jest neurotransmiterem, to powoduje uzależnienie i zwiększa ryzyko wszystkich, bez wyjątku chorób neurologicznych. (6:46) ADHD, stwardnienie rozsiane, stwardnienie zanikowe boczne, choroba Parkinsona. (6:51) To jeszcze nie wszystko.
(6:53) Ten duecik aspartam z glutaminianem przełamują barierę krew-mózg. (7:00) W ich obecności do mózgu mogą wnikać metale ciężkie, rtęć, kadm, ołów, arsen. (7:10) W ich obecności w tkance nerwowej kumuluje się glin.
(7:16) Zwróćcie Państwo uwagę na fakt, że 30 lat temu choroba Alzheimera, a u chorych na Alzheimera stwierdza się duże stężenie glinu, (7:25) nie była tak często diagnozowana. (7:27) Jeżeli w ogóle była diagnozowana, to pojedyncze przypadki były diagnozowane. (7:31) Więc zagrożenie aspartamu wynikające z wszech obecności aspartamu i glutaminianu, (7:40) to jest główna, oczywiście nie jedyna, ale główna najważniejsza przyczyna chorób psychicznych, (7:46) neurologicznych i neurodegeneracyjnych.
(7:49) Co do tego nie ma wątpliwości. (7:50) Podobnie na układ nerwowy działają syntetyczne barwniki, zwłaszcza te koloru niebieskiego. (7:58) Bardzo modne ostatnio są lody i napoje w kolorze niebieskim.
(8:01) W Wielkiej Brytanii wykazano, że dzieci, które spożywały żywność barwioną różnymi barwnikami, (8:09) zachowywały się nienormalnie, miały problemy z koncentracją uwagi, (8:14) z zapamiętywaniem i były nadpobudliwe, hiperaktywne. (8:18) Na wniosek Wielkiej Brytanii w 2019 roku obniżono dopuszczalne dzienne spożycie (8:24) dla kilku tych barwników, dla jednego z nich aż dwudziestokrotnie. (8:30) I mamy potwierdzenie tego, że wyliczone na podstawie LD50 dopuszczalne dzienne pobranie (8:36) to jest ściema, po prostu ściema.
(8:40) Barwnik tartrazyna. (8:41) Nie wiem, czy państwo pamiętacie, lat temu kilkanaście, czterdzieści kilkoro dzieci (8:46) w Miechowie pod Krakowem wylądowało w szpitalu z ciężkim zatruciem. (8:51) Właśnie objawy behawioralne.
(8:54) Domorosły producent słodyczy dobarwił tartrazyną, ładny barwnik, (9:01) taki pomarańczowo-czerwony, jasno-czerwony kolor, koralowy i dzieci się zatruły. (9:08) Ten zakład przestał istnieć, ale tartrazyna w Polsce jest stosowana. (9:13) Wszędzie, gdzie to możliwe.
(9:15) Makarony, płatki kukurydziane, chipsy, napoje, desery, galaretki, (9:19) musztardy, gumy do rzucia również. (9:22) Farmacja stosuje tartrazynę do barwienia syropów, otoczek w pastrykach, (9:26) kapsułek, żeli, balsamów. (9:28) W Szwecji, w Norwegii, w Szwajcarii wycofana jest mniej więcej dziesięć lat temu.
(9:35) Ale nie w Polsce. (9:39) Toksyczne dla mózgu jest połączenie barwników i benzoesanów. (9:45) A benzoesany to są najpowszechniej stosowane konserwanty.
(9:49) Jeżeli ktoś kupuje gotowe suróweczki czy gotowe dania obiadowe, (9:52) no to niestety z benzoesanami ma do czynienia. (9:56) I w obecności benzoesanów witaminy z grupy B są kasowane. (10:01) Ale nie wnikają, to się też przekłada na zwiększone ryzyko chorób neurologicznych, neurodegeneracyjnych.
(10:08) Niezależnie od dodatków, bo tylko tyle w skrócie, działania neurotoksyczne, działania kancerogenne (10:15) wykazują przemysłowe izomery trans obecne w margarynach, w produktach, (10:23) w których margaryny są wykorzystywane, czyli we wszystkich słodyczach, (10:26) zwłaszcza tam, gdzie jest krem batoniki, snickersy, tam gdzie jest krem jest najwięcej tych szkodliwych tłuszczów. (10:33) Również trudne produkty oksydacji nienasyconych kwasów tłuszczowych tak działają, (10:37) więc jeżeli ktoś smaży na olejach, no to ryzykuje własne zdrowie za własne pieniądze. (10:44) Żywność wysoko przetworzona.
Już Państwo macie orientację. (10:50) Białka utlenione, zglikowane wskutek przechowywania długotrwałego. (10:55) Soja genetycznie modyfikowana ze wszystkimi atrakcjami.
(11:01) Transglutaminaza, która powoduje powstawanie niestrawnych wiązań. (11:06) I 10 tysięcy dodatków spożywczych dopuszczonych. (11:10) Około 3 tysiące wykasowano, coś się dzieje, ale tylko 10 tysięcy zostało.
(11:16) W diecie przeciętnego Brytyjczyka żywność wysoko przetworzona stanowi 50\%, (11:21) a w diecie małolatów w Europie 80\%. Nie ma szansy, żeby były zdrowe. (11:26) Po prostu nie ma szans na to.
(11:29) Przez 10 lat monitorowano stan zdrowia prawie 200 tysięcy Brytyjczyków w wieku 40-69 lat (11:36) i wykazano niezaprzeczalny związek między zachorowalnością na nowotwory, (11:44) a spożyciem ultraprzetworzonej żywności. (11:47) Ryzyko raka piersi wzrasta o 16\%, raka płuc o 25\%, ryzyko raka jajnika o 30\%, (11:55) a raka mózgu, uwaga, o 52\%. (11:59) Te wszystkie efekty się po prostu kumulują.
(12:04) W Stanach Zjednoczonych porównywalny dramat. (12:09) Z danych bazy amerykańskiego Centrum Kontroli Prewencji Chorób wynika, (12:13) że w ciągu 20 lat zachorowalność na raka jelita grubego, (12:18) który przez długie lata był traktowany jako choroba seniorów, (12:22) wzrosła w grupie 10-14-latków pięciokrotnie, (12:28) w grupie dzieci, młodzieży 15-19-letniej trzy i półkrotnie (12:35) i w grupie starszej młodzieży 20-24 lata dwukrotnie. (12:40) Rak jelita dzięki temu, że żywność jest jaka jest, (12:47) przestał być uważany za chorobę wyłącznie osób starszych.
(12:51) W Polsce sytuacja nie wygląda wcale lepiej, (12:54) chociaż jest trudniejszy dostęp do danych. (12:56) Z danych Krajowego Rejestru Nowotworów wynika, (12:59) że w 2023 roku odnotowano 181 tys. nowych zachorowań (13:05) i 96 tys.
zgonów z powodu raka. (13:09) Największa dynamika zachorowalności dotyczy młodych ludzi. (13:17) Grupa wiekowa 24-40 największe wskaźniki zachorowalności.
(13:23) W Polsce około 40\% młodzieży 19-24 lata stanowią weganie. (13:30) I teraz warto się zastanowić, (13:32) czy zachorowalność właśnie w tej grupie to przypadek. (13:35) Nie, nie przypadek.
(13:36) Bo jeżeli nie mamy odpowiedniej ilości pełnowartościowego białka (13:42) z aminokwasami siarkowymi na czele, (13:46) to przewód pokarmowy, a konkretnie śluzówka przewodu pokarmowego (13:50) nie może się regenerować. (13:54) A przy przewlekłych stanach zapalnych, już mówiłam, (13:57) dochodzi do rozszczelnienia i dramatu, (14:00) dramatycznej sytuacji ze zdrowiem. (14:02) Dotyczy to wszystkich grup wiekowych, (14:04) ale najczęściej się stwierdza wśród wegetarian i wegan.
(14:12) Dlatego tak promowany jest weganizm. (14:15) No przecież nie bez powodu. (14:18) Nie bez powodu.
(14:19) Ludzie są ogłupiani, dają się ogłupiać. (14:21) To jest najgorsze. (14:23) Słuchajcie, no dokładnie.
(14:27) Można wymówić, że bogaci Polacy, prawda. (14:31) Europejczycy, nie Polacy. (14:33) I młodzi to kupują.
(14:36) Słuchajcie Państwo, te same powody, (14:40) te same przyczyny przekładają się na ryzyko chorób neurologicznych. (14:47) Bo poprzez rozszczelnione jelito, (14:50) do krwiobiegu trafiają niestrawione białka, (14:53) antygeny białkowe, antygeny wytwarzane przez mikroflorę (14:57) i to wszystko zatruwa mózg. (15:00) No zdrowie zaczyna się w jelitach.
(15:02) 80\%, formalnie rzecz biorąc, (15:05) 70\% odporności generujemy w jelitach. (15:08) 20\% odporności zależy od genów, (15:11) a 10\% od opieki medycznej. (15:13) No więc na co stawiać.
(15:15) Ale geny można wyciszyć i można je aktywować. (15:20) Więc z tych 20\%, 10\% można odjąć do zdrowej diety (15:24) i wychodzi, że 80\% odporności generujemy w jelitach. (15:28) Gwałtowny wzrost zachorowalności na choroby psychiczne (15:31) w ciągu ostatnich 20 lat (15:34) to konsekwencja diety zachodniej.
(15:38) Charakterystyczne dla tej diety (15:40) niedobory pełnowartościowego białka, (15:43) niedobory witamin, związków mineralnych (15:45) są przyczyną osłabionych zdolności detoksykacyjnych organizmu, (15:53) a te zdolności są jeszcze świadomie i celowo obniżane. (15:58) Nie wiem, czy państwo wiecie, (16:00) że zalecane spożycie jodu w Polsce (16:02) jest pięciokrotnie mniejsze niż w USA (16:05) i dziesięciokrotnie mniejsze niż w Japonii. (16:10) No bo niedobór jodu to nowotwory piersi, tarczycy, żołądka (16:16) i dzieci z wadami wrodzonymi.
(16:19) Jeżeli dziecko w życiu płodowym nie ma odpowiedniej, (16:22) jeżeli kobieta w ciąży nie ma odpowiedniej ilości jodu, (16:27) to dochodzi do wad wrodzonych. (16:29) To jest świadome i celowe i to jest argument, (16:31) którego się nie da podważyć. (16:34) Ze statystyk Centrum Medycznego NLMED wynika, (16:38) że w 2022 roku liczba wizyt u psychiatry i psychologa (16:41) wzrosła o 30\%, a rok później o 50\%.
(16:46) Wyniki badań kompetencji osób dorosłych z krajów OECD, (16:50) te badania są prowadzone raz na 4 lata, (16:53) więc w badaniach z 2003 roku wykazano, (16:56) że zdrowie psychiczne i zdolności rozumowania wśród Polaków (17:03) wtąpnęły o 30\% w ciągu 12 lat. (17:07) Między 2012 a 2023 rokiem o 30\% zgłupiliśmy, mówiąc wprost. (17:16) A fakt, że jeszcze gorzej wypadli mieszkańcy Chile (17:19) z 31 krajów, to jest raczej wątpliwym pocieszeniem.
(17:24) 11\% nastolatków, czyli ponad 410 tysięcy dzieci i młodzieży (17:30) ma problemy natury psychicznej. (17:32) W 2015 roku samobójstw odnotowano 19, a w 2023 146. (17:39) I tendencja jest rosnąca, nie bez powodu.
(17:44) Ponieważ zdrowie psychiczne kosztuje budżet państwa (17:48) mnóstwo pieniędzy, bo najkrótsze L4 to jest 4 tygodnie. (17:53) A bywa, że delikwent jest leczony 3 miesiące i 6 miesięcy. (17:57) I potem się okazuje, że nie rokuje, idzie na rentę.
(18:00) Więc mam nadzieję, że ten argument powinien rządzących zainteresować. (18:06) Chociaż w tą ekipę rządzącą to absolutnie nie wierzę. (18:09) Nie będę sobie psuć humoru.
(18:11) Więc intelektualna niemoc dzieci i młodzieży, (18:13) podobnie jak to 30\% tąpnięcie dorosłych Polaków nie bierze się znikąd. (18:22) Słuchajcie Państwo. (18:24) Kolejny dowód na to, że nie są to przypadkowe kwestie.
(18:31) Globaliści, nazwisk nie będę wymieniać, (18:34) autorzy tak zwanego wielkiego resetu w 2005 roku w Davos (18:40) postanowili, że mięso, mleko ma z naszej diety zniknąć definitywnie. (18:47) WHO ma już zapisane, że do 50 roku o połowę mniej będziemy tych produktów zjeść. (18:54) Dlaczego? (18:55) No krowy rzekomo wytwarzają metan.
(18:58) Ale eksperci unijni nie wzięli pod uwagę faktu, (19:03) że krowy, które jedzą trawkę, siano czy siano-kiszonki wytwarzają tego metanu. (19:08) Myślę, że pięć, a może i siedem razy mniej. (19:12) Lekarze weterynarze, profesorowie weterynarii twierdzą, że co najmniej pięć razy mniej.
(19:19) Specjaliści od klimatu nie biorą pod uwagę tego, jakie ilości metanu (19:24) uwalniane są do atmosfery przy szczelinowaniu w związku z wydobyciem gazu łupkowego. (19:31) Przepytałam największych ekspertów na spotkaniu Głównego Inspektora Sanitarnego w 2020 roku. (19:37) Żaden nie wiedział ile.
No ale tylko metan od krów jest szkodliwy. (19:41) Więc i mleko, i mięso jest szkodliwe dla biznesu farmaceutycznego. (19:48) Dlaczego? (19:49) W tych produktach występuje wszystko, co służy naszym szarym komórkom.
(19:56) W tych produktach występuje wszystko, co wspomaga nasz układ odpornościowy. (20:03) Krowy na super zbilansowanej monitorowanej komputerowo diecie żyją 4 lata. (20:09) Człowiek na paszy dla zwierząt, bo taką nam wyprodukują.
(20:13) Na końcu prezentacji są trzy filmy i proponuję, żebyście Państwo studiując ten materiał (20:19) zaczęli od tych filmów. To wtedy wszystko będzie jasne i oczywiste. (20:24) Te filmy to jest moje alibi na wypadek, gdyby jakiś TVN-owiec usiłował mnie ścigać.
(20:29) Jednemu się udało, ale nie do końca, bo usiłował zlikwidować firmę, (20:35) która jako jedyna produkuje mleko modyfikowane dla dzieci. (20:38) Efekt był taki, że ja odpuściłam, bo stwierdziłam, że jestem uodporniona (20:43) i na głupotę, i na inwektywy. I po drugie, że z idiotami się nie dyskutuję, (20:48) bo jeżeli dostał moje wystąpienie w sympozjum... (20:53) To co my mamy zrobić jako posłowie, Pani Profesor? (20:56) Ja mam wszystko poukładane, bo mam świetnych sąsiadów prawników.
(21:03) Jeżeli dostał wykaz publikacji, współautorką jest medyczka z linkami, (21:09) co by się nie przepracował przy wyszukiwaniu. (21:12) Jeżeli dostał moją opinię dla Głównego Inspektora Sanitarnego, (21:16) ale on tego w ogóle nie czytał, bo jemu producenci, nawet wiem, (21:20) która firma niemiecka, zapłaciła i napisał. (21:23) Ale efekt był taki, to jest dowód na to, że jeszcze myślący ludzie istnieją.
(21:28) Szef firmy po trzech dniach wykonał telefon, czy nie zmieniłam zdania. (21:33) Ja mówię, nie, z idiotami się nie dyskutuję, szkoda w ogóle czasu. (21:36) A on mówi, to teraz informacja.
(21:39) Magazyny puste, nowy kontrahent dwie tony do Belgii będzie odbierał regularnie. (21:47) Produkujemy na dwie zmiany. (21:48) Firma musiała się wynieść do Czech i na Węgry, (21:50) bo polska biurwokracja od sześciu lat nie robiła nic, tylko niszczyła tę firmę.
(21:58) Tak się traktuje polskich przedsiębiorców, producentów zdrowej żywności. (22:03) To wszystko jest załapówki od koncernów. (22:06) Słuchajcie Państwo, więc przyjęto założenie, że mięso czerwone ma zniknąć (22:11) i pranie mózgów trwa.
(22:13) Ja znalazłam siedemnaście publikacji, w których wykazuje się metodą analizy statystycznej, (22:21) a wiecie doskonale, że jest prawda, półprawda i statystyka, (22:25) że czerwone mięso jest rakotwórcze. (22:28) Pierwsza taka praca powstała oczywiście w USA. (22:31) W 2011 rok wykazano, że jeżeli ktoś regularnie codziennie spożywa 100 gram wołowiny, (22:38) no nie sądzę, żeby ktoś spożywał regularnie codziennie, (22:40) to ryzyko raka rośnie o 17\%.
(22:43) A jeżeli spożywa codziennie, regularnie 50 gram wołowiny przetworzonej, (22:47) to ryzyko raka rośnie o 18\%. (22:50) Więc generalnie w nauce statystyka to jest kwiatek do kożucha. (22:57) Trzeba robić analizy dwiema różnymi metodami i muszą być kompatybilne wyniki.
(23:06) A statystyka służy do podsumowania. (23:09) Trzeba przedstawić dowody, że wołowina jest szkodliwa, (23:14) jak działa, dlaczego, na jakie enzymy, na jakie białka, (23:20) jak się zachowuje w przewodzie pokarmowym. (23:22) Oczywiście tego nie zrobiono.
(23:24) Ale jak się lekarzom na sympozjach tłoczy tą wiedzę, (23:28) nie wiem jak często oni jeżdżą na te sympozy, (23:30) ja podejrzewam, że 3-4 razy w roku, (23:33) to nic dziwnego, że po powrocie z Katowic zostałam paręnaście maili, (23:36) ile prawdy jest w tym, że wołowina jest toksyczna. (23:40) Ale to jest też dowód na to, że zaczynają myśleć. (23:43) Więc w 27 badaniach zrealizowanych w Europie (23:47) tego newsa o rakotwórczych właściwościach wołowiny nie potwierdzono.
(23:52) Dlaczego? (23:54) Wówczas w Europie nie stosowano hormonów w produkcji zwierzęcej. (23:59) To znaczy one formalnie rzecz biorąc są zabronione. (24:02) Ale Unia Europejska nie dysponuje metodami analitycznymi, (24:07) żeby to sprawdzać.
(24:09) A w zakładach, gdzie są wdrożone tak zwane systemy zarządzania jakością, (24:13) nie sprawdza się produkt, tylko przewala się dokumenty. (24:17) Tak to działa. (24:19) Więc macie Państwo dowód na to, (24:23) że pod merką suru jest wszystko przygotowane.
(24:26) Wszystko jest przygotowane. (24:28) Przepisy są. (24:29) Trwałość 23 dni.
(24:31) Transport z Ameryki Południowej do Europy drogą morską (24:35) trwa 11 do 14 dni. (24:38) Jest rezerwa na dystrybucję po sieciach handlowych (24:42) i rezerwa czasowa na sprzedaż tego toksycznego mięsa. (24:47) Więc odpisałam wszystkim bez wyjątku, (24:49) że foliowane, szprycowane, wymieniłam czym, (24:54) foliowane, przechowywane nie wiadomo jak długo, (24:57) bo my nie wiemy jak długo ono było przechowywane, (25:01) jest rzeczywiście trucizną i że trzeba kupować na rynku (25:05) bezpośrednio od producenta.
(25:09) Coś się dzieje, bo czy jestem u córki w Milanówku czy w Olsztynie, (25:12) to na rynku w sobotę nie sposób zaparkować. (25:16) Ale dzięki temu rodzimi producenci odbijają od dna, (25:20) a reszta się truje. (25:22) Więc w hodowli zwierząt w USA, (25:25) w krajach Mercosur stosuje się hormony, (25:29) hormony wzrostu i estradiol, (25:33) którego rakotwórcze właściwości są znane od co najmniej 40 lat.
(25:38) I takie mięso rzeczywiście jest rakotwórcze. (25:41) Te hormony powodują, że w mięsie wzrasta poziom (25:44) insulinozależnego czynnika wzrostu (25:47) i proporcjonalnie do tego czynnika wzrostu (25:50) wzrasta ryzyko nowotworów prostaty i piersi, (25:53) uwaga, odpowiednio 14 i 7-krotnie. (25:59) Amerykańskie niezależne organizacje konsumenckie to sprawdziły.
(26:03) Nie wiem, czy zamieściłam cytowanie, (26:05) ale jeżeli ktoś z państwa będzie potrzebował, (26:07) to są trzy cytowania, trzy dowody z badań naukowych (26:10) na to, że to mięso jest rakotwórcze. (26:13) I teraz uwaga, w badaniach, w których usiłuje się wykazać, (26:17) że czerwone mięso jest przyczyną chorób kardiometabolicznych, (26:20) czyli choroby niedokryjonej serca, że jest rakotwórcze, (26:25) co uznano za przyczynę tych schorzeń? (26:29) Aminokwazy siarkowe i żelazo chemowe. (26:33) Ja nie mam cienia wątpliwości co do tego, (26:35) że to były badania na zamówienie.
(26:39) W jednym przypadku autorzy w stopce pod informacją, (26:42) że deklarują brak konfliktu interesów, (26:45) podali, że finansowane były badania (26:47) przez brytyjski koncern farmaceutyczny. (26:51) Z 17 tych prac 14 autorami byli profesorowie farmacji. (26:57) I tego już nie trzeba komentować.
(27:00) I teraz uwaga, dlaczego aminokwazy siarkowe są dyskredytowane? (27:06) Bo są najważniejsze dla naszych szarych komórek. (27:10) Jeżeli nie mamy w diecie aminokwasów siarkowych, (27:13) to nie możemy wytwarzać glutationu. (27:16) A glutation jest pierwszym najważniejszym antyoksydantem (27:20) dla mózgu, dla układu odpornościowego.
(27:23) W zasadzie działa w całym organizmie. (27:25) Więc aminokwazy siarkowe są dyskredytowane, (27:28) bo szkodzą biznesom. (27:31) Dzięki obecności grupy solohydrylowej (27:34) niestety są one bardzo wrażliwe i bardzo szybko się utleniają.
(27:38) A ich funkcje biologiczne można wymieniać godzinę. (27:41) Najważniejsze determinują strukturę białek. (27:43) Integralność każdej komóreczki (27:46) wchodzą w skład białek strukturalnych, (27:49) hormonów poprzez systemy tiolowe (27:52) zapobiegają oksydacyjnym uszkodzeniom mitochondriów.
(27:56) Łagodzą stres oksydacyjny, (27:58) zwiększają zdolności detoksykacyjne (28:01) na poziomie komórkowym i na poziomie całego organizmu. (28:05) Są niezbędne w endogennej, (28:07) czyli w organizmie, syntezie glutationu, (28:11) który działa detoksykacyjnie, antyoksydacyjnie. (28:15) Są niezbędne w syntezie L-karnityny, kreatyny, (28:19) które regulują metabolizm energetyczny, (28:21) tauryny, która jest neuroprzekaźnikiem (28:23) i metalotioneiny, to jest białko, (28:27) które jest syntetyzowane w organizmie (28:29) wobec różnych zagrożeń, (28:30) na przykład przy zatruciu metalami ciężkimi (28:32) i przy chemioterapii.
(28:34) Również dlatego lekarze onkolodzy (28:35) bardzo tej metalotioneiny nie lubią. (28:39) Przy niedoborach metioniny (28:41) w wątrobie indukowany jest stres oksydacyjny (28:44) ze względu na wyczerpanie zapasów glutationu. (28:48) Co więcej, ograniczone jest wytwarzanie siarkowodoru, (28:51) a siarkowodór pobudza wytwarzanie tlenku azotu, (28:55) który rozszerza naczynia krwionośne.
(28:59) Tlenek azotu reguluje ciśnienie krwi, (29:03) przeciwdziała stanom zapalnym, (29:06) wpływa również na płodność upanów. (29:09) Upani nie, upanów tak. (29:13) Najlepszym źródłem białka, (29:16) które powoduje wytwarzanie dużych ilości tlenku azotu (29:19) są orzechy.
(29:21) Zawierają dużo argininy i arginina właśnie tak działa. (29:27) I wracając do tych 17, (29:29) to znaczy tych badań już jest 40, (29:31) ale ja ich nie analizowałam, (29:33) bo po prostu szkoda na to czasu. (29:35) Więc wracając do tych 17 badań, (29:37) gdzie w 14 przypadkach autorami (29:39) są profesorowi farmacji.
(59:31) Siedemnastu badań, gdzie w czternastu przypadkach autorami są profesorowie farmacji.\\
\% do poniższych czasów dodać 59:36 \%\\
(0:00) Dziwnym trafem, zbierając dane od pacjentów, nie uwzględniono absolutnie żadnych czynników ryzyka oprócz tego czerwonego mięsa. (0:17) Nie uwzględniono otyłości cukrzycy, nadciśnienia, palenia papierosła.
(0:24) Przecież wiadomo, że otyłość i cukrzyca zmieniają metabolizm radykalnie (0:28) i zwiększają 500 razy ryzyko wszystkich innych chorób. (0:32) Nie uwzględniono infekcji helikobakter pylori, również niskiej aktywności fizycznej. (0:39) W rzeczywistości przyczyn choroby niedokrwiennej serca było na pewno co najmniej kilka, jeśli nie kilkanaście.
(0:47) Smażenie na olejach, grillowanie, mówiłam o tym. (0:51) Produkty oksydacji lipidów, które utleniają białka, utleniają witaminy. (0:57) Przy długotrwałym magazynowaniu niedobory witamin z grupy B, a to się przekłada na hiperhomocysteinę.
(1:04) Każdy z nas wytwarza homocysteinę, bo nie ma innych szlaków metabolicznych. (1:11) Jeżeli natomiast mamy odpowiedni poziom B6, B9, B12, to ta homocysteina jest redukowana do cysteiny, (1:20) z której powstaje glutacjon, tauryna, metalotioneina i nie tylko. (1:26) Więc przy niedoborach witamin z grupy B, zero szans na zdrowe szare komórki.
(1:32) Dlaczego jako wyłączną przyczynę reakcji wolnorodnikowych autorzy badań potraktowali żelazo chemowe? (1:39) No bo tak im pasowało. (1:40) Wyniki badań kohortowych są interpretowane tendencyjnie pół biedy, ale są interpretowane w całkowitej sprzeczności z biochemią. (1:51) Czyżby badania na zamówienie, to znaczy ja nie mam wątpliwości, (1:56) za chorowalność i na chorobę niedokryjoną serca i na nowotwory, (2:02) koreluje w statystycznie istotnym stopniu tylko ze spożyciem mięsa przetworzonego, strach z glutaminazą, soją i dodatkami.
(2:13) Nie koreluje ze spożyciem mięsa nieprzetworzonego, no ale kto by się w przypadku farmacji takimi detalami przejmował. (2:21) Dla nich to jest nieistotne. (2:22) Proszę następny, bo ja już chyba mam niewiele czasu.
(2:26) Jakie są zdrowotne konsekwencje rezygnacji ze spożycia białka zwierzęcego? (2:32) To i poprzedni slajd o aminokwasach siarkowych to jest clue dzisiejszego wykładu. (2:38) Niedobory aminokwasów egzogennych, niedobory żelaza, cynku i witaminy B12. (2:47) Przy tych niedoborach nie ma szans ani na prawidłowy rozwój dziecka, ani na zachowanie zdrowia.
(2:56) Po prostu nie ma szans. (2:57) Żelazo, witamina B12 to jest dotlenienie tkanek z mózgiem na czele. (3:02) Cynk to hormony wzrostu, to również hormony płciowe u panów, a nawet łysinie zapobiega cynk.
(3:11) U dzieci, które mają we wczesnym dzieciństwie niedobory cynku obserwuje się krótkie nóżki. (3:17) W Kenii wykonano trwające 3,5 roku badania w grupie dzieci, które jadły tradycyjnie dietę z mięsem. (3:26) Wykazano zdecydowanie lepszy rozwój fizyczny, a zwłaszcza psychiczny w porównaniu do grupy, która była na diecie wegen.
(3:34) To były dzieci w wieku wczesnoszkolnym, 8-11 letnim. (3:39) Niedobory pełnowartościowego białka to osłabiona osłona antyoksydacyjna i ryzyko stanów zapalnych. (3:45) Tu aminokłasy siarkowe się kłaniają.
(3:48) Niedotrenienie tkanek, zaburzenia w funkcjonowaniu serca mózgu, zaburzenia równowagi hormonalnej, (3:55) obniżony poziom hormonów tarczycych, hormonów wzrostu oraz insulinozależnego czynnika wzrostu plus problemy z płodnością. (4:04) Ograniczona zdolność syntezy neurotransmiterów, co zwiększa ryzyko zaburzeń emocjonalnych, bezcenności, depresji, neurodegeneracji. (4:16) Może dochodzić do zaniku mięśni, pogorszenia sprawności fizycznej.
(4:21) Wzrasta wówczas ryzyko złamań nie tylko u seniorów, no i większe ryzyko śmiertelności z najprzeróżniejszych przyczyn. (4:31) Słuchajcie Państwo, mięso to jest dieta dla mózgu. (4:35) Są już dowody naukowe od niedawna na to, że neandertalczyk był mięsożercą.
(4:43) Jego objętość czaszki w ciągu 42-44 tysięcy lat zwiększyła się trzykrotnie kosztem szczęki. (4:52) Dowody w postaci neandertalczyka wykopanego gdzieś na Syberii i we Francji wykazały, że w układzie kostnym u tych osobników występuje bardzo duże stężenie cynku. (5:11) A biodostępny cynk, jak już wiecie zapewne, jest obecny w produktach pochodzenia zwierzęcego.
(5:18) Przez trzy i pół roku nie wiedzieli badacze francuscy, bo oni te badania zrobili, jak to odnieść do roślinożerców i mięsożerców. (5:28) I jakiś doktorant wpadł na pomysł znaleźć w muzeach szkielety zwierząt sprzed 40 tysięcy lat. (5:38) Pół roku namierzali, drugie pół roku załatwiali pozwolenie, no bo trzeba było jeden gram nawiercić i pobrać.
(5:47) No i bardzo szybko badania doktorant zrobił w trzy miesiące i okazało się, że poziom cynków w szkielecie neandertalczyka był taki sam jak u rysia i u wilka. (6:00) I cztery razy większy niż u królika i kozicy, niż u roślinożerców. (6:05) I to jest dowód nie do podważenia.
(6:08) Natomiast to, co się aktualnie dzieje zmierza dokładnie w odwrotnym kierunku. (6:13) W Stanach Zjednoczonych i w Holandii od lat siedemdziesiątych wszyscy poborowi są badani pod kątem wskaźnika IQ. (6:24) I ta inteligencja rok w rok maleje.
(6:29) Czasami o 0,7\%, czasami o 0,4\%. (6:34) W każdym razie po pięćdziesięciu latach wskaźnik IQ, no fakt, że do wojska nie idą ludzie najlepiej wykształceni i najlepiej odżywiający się, ale nie ma innych danych. (6:46) Więc w ciągu pięćdziesięciu lat ten wskaźnik IQ zmalał o 26\%.
(6:51) Po ujawnieniu tych wyników profesor Stanisław Berger z SGGW w Warszawie zrobił eksperyment na szczurach. (6:59) Grupa zwierząt dostawała tłuszcze zwierzęce. (7:02) Grupa zwierząt dostawała tłuszcze roślinne.
(7:05) I zgadnijcie Państwo jak duże różnice były w objętości czaszki u zwierząt. (7:12) Badano tam ich zdolności manualne, zachowania w labiryncie, poszukiwanie jedzenia, ale najlepszym wskaźnikiem jest waga mózgu. (7:20) Więc waga mózgów szczurków żywionych olejami była o 28\% mniejsza niż tych, które jadły smalczyk, masło, tłuste sery.
(7:33) Wszystkie pozostałe składniki w diecie były takie same. (7:36) Więc rosnąca zachorowalność, jak już wspominałam chyba trzy razy, to nie jest żaden przypadek. (7:44) To jest świadome i celowe generowanie chorób.
(7:49) Generowanie chorób i depopulacja przy okazji. (7:53) Dziennikarzowi z TVN-u, który mnie obsmarował za spiskową teorię dziejów, wysłałam oprócz materiałów, o których mówiłam, jeden wykres. (8:01) Nie mam go na pendrive.
(8:02) Stopa zysku 200 największych przedsiębiorstw w USA w procentach, w tym farmacji i stopa zysku dla samej farmacji. (8:15) I uwaga, w 2010 roku stopa zysku dla farmacji była 5,5 razy większa niż dla tych największych firm z farmacją włącznie. (8:25) Po COVID-zie to może być 10 albo i więcej, nie wiem, nie dokopałam się do danych.
(8:32) A nie mam już zajęć ze studentami, bo zawsze ich wykorzystywałam do archiwizacji danych, bardzo to lubili robić. (8:39) Ale może kogoś w to wpuszczę, w ten temat. (8:42) Obciążenie chorobami wcale nie wynika z nadmiaru białka w diecie, co usiłują nam wmawiać lekarze.
(8:48) Wcale nie wynika z konsumpcji czerwonego mięsa pod warunkiem, że jest to mięso, a nie produkt mięsopodobny. (8:56) Nie wynika z całą pewnością z nadmiaru aminokwasów siarkowych. (9:00) Jest konsekwencją dużego spożycia żywności przetworzonej, która zawiera utlenione, glikowane, nitrozowane białka.
(9:08) Białka, które nie są trawione w naszym organizmie. (9:11) Jednocześnie ta dieta charakteryzuje się około dziesięciokrotnym niedoborem antyoksydantów, pierwiastków śladowych i witamin z grupy B. (9:23) Nie ma szans, żeby być na niej zdrowym. (9:26) Byle jakość żywności przetworzonej to nie przypadek.
(9:30) Najlepszym tego dowodem jest stosowanie w przetwórstwie żywności soli kuchennej, (9:37) tej wzbogaconej, w jodan potasu i żelazu cianek potasu, który pełni funkcję antyzbrylacza. (9:47) Jodan działa prooksydacyjnie, intensyfikuje procesy utleniania żywności. (9:53) W próbkach mięsa z solą jodowaną i jodanym spadek zawartości biodostępnej, lizyny i metioniny, to jest aminokwas siarkowy, (10:01) jest średnio o 20\% większy w porównaniu do próbek z solą niejodowaną bądź jodowaną jodkiem.
(10:10) Stosowanie jodanu i żela z ocianku to dowód na świadome, celowe obniżanie wartości odżywczej i generowanie chorób metabolicznych, (10:18) które od połowy lat dziewięćdziesiątych traktowane są jako nieuleczalne. (10:24) W połowie lat osiemdziewięćdziesiątych do połowy lat dziewięćdziesiątych cukrzyca była uleczalna. (10:31) Obturacyjna choroba płuc była uleczalna.
(10:34) Zespół metaboliczny był uleczalny, a teraz wszystko jest nieuleczalne, bo tak się branża farmaceutyczna umówiła z lekarzami. (10:44) Dlatego lekarze, którzy usiłują likwidować przyczynę choroby poprzez właściwą podaż białka, witamin, mikroelementów są dyskryminowani, bo psują biznes. (10:57) Przecież wyleczony pacjent to strata dla farmacji.
(11:01) Ma się leczyć do końca życia i jeszcze najlepiej po śmierci parę dni. (11:06) Więc słuchajcie Państwo, wbrew temu, że dopuszczono zgodnie z prawem, tak zwanym prawem żywnościowym, sól wzbogaconą w jodan i żelazo cianek, osoba, która podpisała te dokumenty powinna trafić do kryminału. (11:24) To jest opinia moich sąsiadów prawników.
(11:27) Więc jodan generuje stany zapalne jelita oraz choroby tarczycy. (11:35) Choroby tarczycy to jest obniżone wytwarzanie energii życiowej. (11:41) To się wszystko kumuluje.
(11:42) Żelazo cianek potasu powoduje niedotlenienie tkanek. (11:46) Żelazo cianek blokuje w 20\%, a nawet w 100\% aktywność 20 różnych enzymów z grupy oksydoreduktaz, dzięki czemu tlen przenoszony przez czerwone krwinki nie może być odebrany do tkanek, no bo nie ma enzymów, które to wychwytują. (12:04) A najbardziej wrażliwy na niedotlenienie jest jaki organ w naszym organizmie? (12:11) Mózg.
Dokładnie mózg. (12:15) Słuchajcie Państwo, ja dzięki uprzejmości aktualnego dyrektora Narodowego Instytutu Zdrowia Publicznego, bo z dyrektorem Iżędżetu miałam na pieńku przez 30 lat, stale i niezmiennie. (12:29) Jak byłam w komisji doktorskiej, to nie przyjeżdżał w charakterze recenzenta, tylko przysyłał kwity i odwrotnie.
(12:37) Po prostu alergia. (12:40) Alergia na przekręt. (12:42) Bo wszystkie granty unijne są, ile siana dodać, ile słomy, ile skrobi, ile błonnika, a jaki jest wpływ na wartość żywieniową 0,0,3.
(12:56) Trzy projekty z mojej katedry dotyczące właśnie tych zmian w zawartości witamin i w poziomie antyoksydantów. (13:06) Jeden prosty test o wszystkim mówi. (13:09) Wyautowane.
(13:12) Bo ci, co rządzą i rozdzielają kasę, no są niestety na smyczy nazywając rzeczy prawdziwe. (13:18) Dzięki uprzejmości dyrektora Instytutu Zdrowia Publicznego jestem w posiadaniu obu ekspertyz, na podstawie których wydano w 2010 roku pozwolenie na stosowanie jodanu. (13:33) I tam nie ma nawet słowa o tym, że jodan jest bezpieczny.
(13:37) W ekspertyzie jest napisane, że może powodować stany zapalne, ale nie ma na to dowodów in vivo. (13:47) A przecież formalnie rzecz biorąc w prawie żywnościowym jest zasada ostrożności, więc jeżeli może powodować, to nie powinien być dopuszczony, a został dopuszczony. (14:00) I mniej więcej pięć lat temu dodano do tego żelazo cianek i ten duet tych toksyn powoduje problemy w jelitach, gdzie generowana jest odporność.
(14:13) Chore jelita równa się chory umysł. (14:15) Więc nie ma cudów. (14:17) Ja dziękuję za uwagę.
(14:18) Jeżeli przekroczyłam czas, to sorry. (14:20) W 2022 roku zrobiłam akcję z mlekozastępczymi preparatami do żywienia niemowląt. (14:30) Przedstawiłam argumenty, chyba 60 pozycji literatury, głównie anglojęzycznej.
(14:37) Przedstawiłam swoje własne badania realizowane wspólnie z Akademią Medyczną w Warszawie, trzy publikacje recenzowane. (14:43) Nawet byłam zdziwiona, że opublikowali. (14:46) Oczywiście grantu z Polskiej Akademii Nauk nie dostałyśmy raz, nie dostałyśmy drugi raz, nie dostałyśmy trzeci raz.
(14:53) Ale Krajowy Związek Współdzielni Mleczarskich stwierdził, że kupią nam odczynniki. (14:57) No i wystarczyło. (14:58) Honorariów nie było, ale odczynniki, warte chyba 26 czy 28 tysięcy, wykorzystałyśmy.
(15:07) I nawet miałam cień nadziei, że coś się stanie. (15:11) Dostałam pismo, że z prokuratury jednej przekazano do drugiej, po trzech tygodniach z drugiej do trzeciej, potem z trzeciej do czwartej to było na Pradze. (15:21) Potem znowu gdzieś wróciło, to trwało pół roku.
(15:24) Informacja spór kompetencyjny i przed samą Wigilią dostaję pismo od prokuratora, że skoro stawiam zarzuty, to mam się ustosunkować do konkretnych pytań. (15:36) Mój mąż jak to zobaczył, mówi słuchaj prawnik tego nie pisał. (15:40) Albo lekarz, albo technolog żywności, bo były wyjątkowo merytoryczne.
(15:46) Robiłam akurat recenzję książki, a zawsze wydawnictwa wszystko robią naprzód wczoraj, więc zostawiłam tą recenzję książki. (15:53) Ustosunkowałam się do pytań, wysłałam po czterech czy pięciu dniach, żeby przed świętami to doszło. (15:59) W styczniu zostałam zawezwana na komendę policji, co by złożyć zeznania.
(16:06) No i zeznałam, co jest prawdą. (16:08) Kto dopuścił produkt na rynek? (16:12) Dlaczego nie interweniowano w 2006 roku, kiedy było masowe zatrucie dzieci? (16:16) 13-14 dzieci się zatruło. (16:19) Producent preparatu twierdził, że matki nie zachowały higieny i że to jest infekcja bakteryjna.
(16:25) Większość mam okazało się były lekarkami, więc zrobiły posiewy i dowody w postaci badań z sanepidu, że nie było infekcji bakteryjnej. (16:36) Ja byłam wtedy na sympozjum chyba w Krakowie albo w Tarnowie. (16:39) Padło pytanie, czy da się to udowodnić.
(16:43) Profesor, notabene psychiatri mówi, to ma pani fucha. (16:46) Ja nie mam fuchy, bo ja nie mam sprzętu, żeby to oznaczyć. (16:48) Więc nie zainteresowano się w ogóle jak, dlaczego kobiety zostały same z problemem.
(16:57) Dzięki mojej opinii, dzięki wynikom zrobionym na SGGW w Warszawie miały argumenty i produkt zniknął z rynku. (17:06) Na dwa lata. (17:08) A potem został wprowadzony ponownie.
(17:11) I prokurator nie widział w tym nic zdrożnego. (17:15) Że produkt nie był przebadany przy powtórnym wprowadzeniu na rynek. (17:20) Że w ogóle nie jest badany.
(17:22) Ale prokurator zrobił jedną dobrą rzecz. (17:25) Od Ministerstwa Zdrowia wycyganił sposób oceny tych preparatów. (17:30) Od sasa do lasa na niby, aby, aby.
(17:33) W tym zrobimy metale ciężkie, a tu kwas erukowy. (17:36) A przecież oleje wysokoerukowe zniknęły z rynku w latach 70-tych. (17:41) Więc nie wiem po co kwas erukowy oznaczali.
(17:44) Cały nadzór jest na niby, na aby, aby. (17:48) Rok później zrobiłam, więc prokurator w końcu po tych wszystkich perypetiach przysłał mi zresztą eleganckie pismo. (17:56) Że to jest wszystko zgodne z prawem żywnościowym.
(17:59) No bo jest. No bo jest. (18:02) Rok później zrobiłam akcję z solą.
(18:04) Dwa zakłady odpisały, że nie jodują w ogóle. (18:09) Kilka zakładów, że jodują jodkiem, ale wiem, że te zakłady teraz jodują jodanem. (18:15) Bo jest to obowiązkowe.
(18:17) W 2001 roku wprowadzono jodan albo jodek do wyboru. (18:23) A w 2000 wprowadzono obligatoryjnie jodan, no po to, żebyśmy zdrowi nie byli. (18:28) Więc jodują.
(18:31) Jedynie z Kłodawa nie joduje w ogóle. (18:35) Ale dziwnym trafem w Czechach, na Węgrzech, na Słowacji joduje się jodkiem niejodanym. (18:42) Czyli unijne prawo nie jest nadrzędne.
(18:45) Bo nie jest nadrzędne. (18:47) I w akcie desperacji trochę ochłonęłam parę miesięcy po stresie. (18:54) W ubiegłym roku, dokładnie chyba rok temu, w marcu albo kwietniu, do pięciu posłów wysłałam informację o wszystkich zagrożeniach.
(19:02) To w tekst 12 stron. (19:05) I jedynie pan poseł Grzegorz Braun zareagował. (19:11) Chociaż nie jest biochemikiem, nie jest lekarzem, ale potrafi myśleć w kategoriach dobra wspólnego.
(19:21) Pozostali posłowie, cztery osoby, a wybierałam takie osoby, które mają żony lekarki. (19:28) Pan poseł Roman Fritz nie był posłem. (19:31) Ja do pana posła Fritza nie wysyłałam tego.
(19:34) A kiedy to było wysyłane mniej więcej? (19:36) W marcu ubiegłego roku albo w lutym. (19:41) Więc reakcji ze strony urzędników nie ma żadnej. (19:46) Jest zainteresowanie.
(19:47) Bo dyrektor Instytutu Zdrowia Publicznego nie musiał mi tych ekspertyzm przysyłać. (19:53) Bo nie musiał. (19:54) Zrobił to z własnej nieprzymuszonej woli.
(19:56) Może dlatego, że równie jak ja nie lubił poprzedniego dyrektora. (20:00) Nie wiem. (20:01) W każdym razie nie musiał mi tego przysyłać.
(20:04) Są dowody na to, że jest to przekręt. (20:07) Także jeżeli jakiekolwiek akcje byłyby związane z bezpieczeństwem, raczej z zagrożeniem naszego bezpieczeństwa, (20:15) to UOKiK powinien pobrać próbki i wykazać w badaniach ile soi, ile mięsa jest w produktach. (20:28) UOKiK jest jedyną agendą rządową, która zawsze w sytuacjach podbramkowych reaguje.
(20:35) Trzykrotnie reagowała. (20:37) Raz usunęli z TVN-u reklamę mleka o unikalnej formule, (20:43) które to mleko zabierało 2,2 białek mleka plus oleje i chemia. (20:50) Ale zaraz potem się pojawiły, chyba nawet wyższe niż ta sala, pionowe banery w supermarketach.
(20:57) Mleko o unikalnej formule na półkach z ekożywnością. (21:02) I na to już nie było bata. (21:05) Już tutaj UOKiK nie mógł, bo sieci handlowe mają swoje zwyczaje i swoje prawa.
(21:10) UOKiK zrobił akcję, chyba trzy lata temu, z usunięciem mleko zastępczych preparatów, (21:19) bo na etykietach było napisane, że produkt na bazie mleka. (21:23) No to wyliczyłam, ile tego mleka zawiera. (21:26) No i wyszło, że zawiera 2\% albo 3\%, a reszta to nie jest mleko.
(21:30) No ale skończyło się niestety tylko na tym, że wycofano w całej Polsce ze sklepów również Rossmana. (21:41) Obraza majestatu była niesamowita. (21:44) Wycofano i zmieniono etykiety.
(21:46) I główny inspektor sanitarny dopuścił to. (21:49) Więc w pierwszej kolejności należałoby zlikwidować Sanepid i GIS. (21:57) Bo to są rządowe agendy, które pilnują interesów koncernów.
(22:07) No ja 5-7 razy się przekonałam, że z nimi w ogóle się nie da dyskutować. (22:13) A pan minister Jan Krzysztof Ardanowski również doświadczył tego samego. (22:18) Usiłował utworzyć na wzór Francji, Niemiec czy Węgier czy Czech (22:22) jeden Urząd do Spraw Bezpieczeństwa Żywności.
(22:25) Ale byłby wtedy jeden dyrektor, dwóch czy trzech zastępców. (22:29) Ale byłyby środki na badania żywności. (22:33) A możliwości są, bo ze środków unijnych wszystkie Sanepidy, GISy mają taki sprzęt, (22:38) że nie jedna uczelnia może pozazdrościć.
(22:40) Ten sprzęt stoi i rdzewieje. (22:43) I kto był oponentem? (22:45) Dzięki komu pomysł niestety się nie ziścił? (22:49) Sanepid i GIS. (22:50) Sanepid.
\end{document}